\documentclass[11pt,a4paper]{article}

\usepackage[margin=25mm]{geometry}
\usepackage[T1]{fontenc}
\usepackage{lmodern}
\usepackage{amsmath,amssymb,amsthm}

\theoremstyle{definition}
\newtheorem{definition}{Definition}
\newtheorem{problem}{Problem}

\setlength{\parindent}{0pt}
\setlength{\parskip}{6pt}

\begin{document}

\begin{center}
  {\Large\bfseries Davenport--Schinzel Sequences: Growing Order Below the Critical Scale}
\end{center}


\section{Definition}
\begin{definition}[Order-$s$ Davenport--Schinzel sequence]
Let $n,s\geq 1$ and $[n]=\{1,\ldots,n\}$. A sequence
$U=(u_1,\ldots,u_m)$ over $[n]$ is an order-$s$ Davenport--Schinzel sequence if
\begin{enumerate}
  \item $u_i\neq u_{i+1}$ for every $i$; and
  \item for every two distinct symbols $a,b$, the sequence $U$ contains no
  alternating subsequence
  \[
    a,b,a,b,\ldots
  \]
  of length $s+2$.
\end{enumerate}
We write $U\in\mathrm{DS}(n,s)$.
\end{definition}

\begin{definition}[Extremal function]
\begin{equation}
  \lambda_s(n)
  =\max\bigl\{|U|:U\in\mathrm{DS}(n,s)\bigr\}.
\end{equation}
\end{definition}

For all $n$ and $s$, the pigeonhole argument gives
\begin{equation}
  \lambda_s(n)\leq\binom{n}{2}(s+1).
\end{equation}

\section{Current best result}
The first values are exact:
\begin{equation}\label{eq:small-orders}
 \lambda_1(n)=n,\qquad \lambda_2(n)=2n-1,\qquad
 \lambda_3(n)=2n\alpha(n)+O(n),
\end{equation}
where $\alpha$ denotes the inverse-Ackermann function.  The fixed-order problem is
now essentially settled.  In particular,
\begin{equation}\label{eq:fixed-order}
 \lambda_4(n)=\Theta\!\bigl(n2^{\alpha(n)}\bigr),\qquad
 \lambda_5(n)=\Theta\!\bigl(n\alpha(n)2^{\alpha(n)}\bigr),
\end{equation}

The known fixed-order
constructions remain effective up to $s=O(\alpha(n))$.

On the other hand,
\begin{equation}
  \lambda_s(n)\sim\binom{n}{2}s
\end{equation}
is known when $s/\sqrt n\to\infty$. There are no matching asymptotic bounds
between these two ranges.

\section{Problem formulation}
Determine the asymptotic behavior of $\lambda_s(n)$ when
\begin{equation}
  s=s(n)\longrightarrow\infty,
  \qquad
  s=o(\sqrt n),
  \qquad
  s=\omega(\alpha(n)).
\end{equation}

\end{document}
